\SourceLeaf{226}%
\noindent Von \(A\) gehen bei jedem der beiden Großkreise zwei Zweige aus. Dort entstehen also zwei Winkel, die einander zu \(180^0\) ergänzen. Unsere Untersuchung zeigt aber, dass wir die Zweige wählen müssen, deren Richtungen mit dem Durchlaufsinn von \(L\) nach \(L'\) und von \(L''\) nach \(L'''\) übereinstimmen. Damit ist zugleich klar, dass wir zwischen den beiden Schnittpunkten der Großkreise frei wählen können. Anstelle des Winkels \(A\) kann man auch den Bogen zwischen den Polen der Großkreise verwenden, zu denen die Bögen \(LL',L''L'''\) gehören. Offenbar sind dabei Pole zu wählen, die bezüglich der beiden Bögen gleich liegen, also etwa beide rechts oder beide links, wenn wir von \(L\) nach \(L'\) und von \(L''\) nach \(L'''\) gehen.

VII. Seien \(L,L',L''\) drei Punkte auf der Kugeloberfläche. Der Kürze halber setzen wir
\[
\begin{aligned}
\cos(1)L&=x,&\cos(2)L&=y,&\cos(3)L&=z\\
\cos(1)L'&=x',&\cos(2)L'&=y',&\cos(3)L'&=z'\\
\cos(1)L''&=x'',&\cos(2)L''&=y'',&\cos(3)L''&=z''
\end{aligned}
\]
\noindent sowie
\[
xy'z''+x'y''z+x''yz'-xy''z'-x'yz''-x''y'z=\Delta
\]
\noindent Es bezeichne \(\lambda\) den Pol des Großkreises, zu dem der Bogen \(LL'\) gehört, und zwar denjenigen, der zu diesem Bogen genauso liegt wie der Punkt \((1)\) zum Bogen \((2)(3)\). Nach dem vorangegangenen Satz gilt \(yz'-y'z=\cos(1)\lambda\mathord{.}\sin(2)(3)\mathord{.}\sin LL'\) oder wegen \((2)(3)=90^0\)
\[
\begin{aligned}
yz'-y'z&=\cos(1)\lambda\mathord{.}\sin LL',\ \text{und ebenso}\\
zx'-z'x&=\cos(2)\lambda\mathord{.}\sin LL'\\
xy'-x'y&=\cos(3)\lambda\mathord{.}\sin LL'
\end{aligned}
\]
\noindent Multiplizieren wir diese Gleichungen der Reihe nach mit \(x'',y'',z''\) und addieren sie, so erhalten wir mit Hilfe des zweiten in \(V\) genannten Satzes
\[
\Delta=\cos\lambda L''\mathord{.}\sin LL'
\]
\noindent Nun sind drei Fälle zu unterscheiden. Im \textit{ersten} Fall liegt \(L''\) auf demselben Großkreis, zu dem der Bogen \(LL'\) gehört. Dann gilt \(\lambda L''=90^0\) und damit \(\Delta=0\). Liegt \(L''\) dagegen außerhalb dieses Großkreises, so tritt der \textit{zweite} Fall ein, wenn der Punkt auf derselben Seite wie \(\lambda\) liegt, und der \textit{dritte}, wenn er auf der entgegengesetzten Seite liegt. In diesen beiden Fällen bilden die Punkte \(L,L',L''\)
